\chapter{Nebenläufige Programmierung: Threads \tagLK}
\renewcommand{\chaptertitle}{Nebenläufige Programmierung: Threads}

\lehead[]{\normalfont\sffamily\hspace*{-2.00cm}\textcolor{white}{\colorbox{lightblue}{\makebox[1.60cm][r]{\thechapter}}}\hspace{0.17cm}\textcolor{lightblue}{\chaptertitle}}
\rohead[]{\textcolor{lightblue}{\chaptertitle}\normalfont\sffamily\hspace*{0.17cm}\textcolor{white}{\colorbox{lightblue}{\makebox[1.60cm][l]{\thechapter}}}\hspace{-2.00cm}}
%\chead[]{}
\rehead[]{\textcolor{lightblue}{AvHG, Inf, My}}
\lohead[]{\textcolor{lightblue}{AvHG, Inf, My}}

\lstset{style=myJava}

\section{Wofür benötigt man Threads?}

Threads ermöglichen es, ein Programm in mehrere Teile zu unterteilen, die
unabhängig voneinander parallel ausgeführt werden. Jeder Thread (engl.\ für
Faden) arbeitet eigenständig in seinem eigenen Tempo. Über Variablen können die
Threads Informationen untereinander austauschen. Wenn ein Programm
beispielsweise eine lange Berechnung durchführen muss, die mehrere Minuten
dauert, so wäre es in einem \glqq normalen\grqq\ Programm während dieser Zeit
nicht möglich, auf Benutzereingaben zu reagieren. Wenn man die Berechnung dagegen
in einem zweiten Thread durchführt, kann sich der Haupt-Thread weiterhin um
Benutzereingaben kümmern. Weitere Beispiele für den Einsatz von Threads:

\begin{compactitem}
\item Wenn ein Text gedruckt wird, ist dies eine Aktion, die relativ lange
dauert. Wird der Druckvorgang in einem eigenen Thread durchgeführt, so steht der
Haupt-Thread dem Benutzer während des Druckvorgangs für weitere Arbeitsaufträge
zur Verfügung.

\item Während auf Eingaben des Benutzers gewartet wird, sollen im Hintergrund
aufwändige Berechnungen fortgeführt werden.
\end{compactitem}

Jede von \myClass{Thread} abgeleitete Klasse muss die Methode \textbf{\texttt{run()}}
überschreiben. Diese wird automatisch ausgeführt, sobald der Thread über die
geerbte (und bei Bedarf ebenfalls überschriebene) Methode
\textbf{\texttt{start()}} aus der Basisklasse \myClass{Thread} gestartet
wird.

Der Methode \textbf{\texttt{start()}} kommt dabei eine besondere Aufgabe zu:
Sie sorgt dafür, dass der neue Thread aus Sicht des Betriebssystems als
eigenständiger Prozess läuft. Nur so laufen die verschiedenen Threads
tatsächlich parallel nebeneinander her. Würde man \lstinline|run()| hingegen
direkt aufrufen (was technisch gesehen durchaus möglich ist), würde diese
Methode -- wie alle anderen Methoden auch -- ganz normal ausgeführt. Und der
nächste Programmschritt würde erst dann bearbeitet, wenn \lstinline|run()| seine
Arbeit beendet hat.

Obwohl es in der Basisklasse \myClass{Thread} eine Methode
\textbf{\texttt{stop()}} gibt, sollte man diese nicht verwenden, um einen
Thread zu beenden! Vielmehr sollte man über eine boolesche Variable, die den
Zustand (läuft oder läuft nicht) des Threads festhält, steuern, ob die Methode
\lstinline|run()| weiter ausgeführt werden soll. Damit alle Threads Änderungen
an einer solchen gemeinsamen Zustandsvariable sofort bemerken, kennzeichnet man
sie zusätzlich mit dem Schlüsselwort \textbf{\texttt{volatile}} (engl.\
\emph{flüchtig}). Ohne \lstinline|volatile| dürfte ein Thread aus
Geschwindigkeitsgründen einen veralteten Wert der Variable in seinem lokalen
CPU-Cache behalten, anstatt ihn frisch aus dem gemeinsamen Hauptspeicher (RAM)
zu lesen. Der Thread als Ganzes wird beendet, sobald die Methode
\lstinline|run()| beendet ist. Siehe dazu das letzte Beispiel in diesem Kapitel.

\begin{center}
\begin{minipage}{1.0\textwidth}
\begin{framed}
\textbf{Merksatz:} \texttt{run()} muss, \texttt{start()} kann
überschrieben werden. Und \texttt{stop()} sollte man nicht verwenden --
stattdessen beendet man Threads sauber über eine mit \texttt{volatile}
gekennzeichnete Zustandsvariable.\\[1mm]
\emph{Moderne Kurzform:} Wenn eine Klasse nicht von \myClass{Thread} erben kann
oder soll, lässt sich ein Thread auch direkt mit einem Lambda-Ausdruck starten:\\
\texttt{new Thread(() -> \{ /* Code für run() */ \}).start();}
\end{framed}
\end{minipage}
\end{center}
